\documentclass[11pt]{article}
\usepackage[margin=1in]{geometry}
\usepackage{booktabs}
\usepackage{longtable}
\usepackage{tabularx}
\usepackage{array}
\usepackage{amsmath}
\usepackage{hyperref}
\usepackage{caption}

\captionsetup{font=small}

\title{CLARITY SemEval 2026 Submission Report\\
Experiment Forensics Across Granite, Qwen, DeBERTa, RoBERTa, and Ensemble Runs}
\author{Workspace: \texttt{/Users/andrearachetta/Desktop/CLARITY-SemEval-2026}}
\date{February 14, 2026}

\begin{document}
\maketitle

\begin{abstract}
This report documents the current CLARITY-SemEval-2026 experimentation workspace with a practical focus on (i) model artifacts and metrics, (ii) repository provenance and git state, and (iii) submission-readiness of generated files. The analysis covers Granite, Qwen, DeBERTa, RoBERTa, and ensemble-related assets found in notebooks, JSON files, trainer checkpoints, and folder structures. The repository behaves as an experiment staging area: most model work is currently untracked by git. Consequently, this report emphasizes file-level evidence and explicit paths.
\end{abstract}

\section{Scope and Evidence Collection}
\textbf{Target requested by user.} Produce a detailed SemEval-style report with comprehensive tables, focusing on git changes and experiment/model artifacts (especially Granite-related submission workflow), including results recoverable from notebooks, JSON files, and model folders.

\textbf{Important naming note.} No folder or file named \texttt{GraniteClaritySubmission} was found. Granite submission logic appears in:
\begin{itemize}
    \item \path{GraniteClarityEvaluation.ipynb}
    \item \path{GraniteClarityEvaluation copy.ipynb}
    \item \path{GraniteClarityEvaluation copy 2.ipynb}
\end{itemize}

\section{Repository Provenance and Git State}
\begin{table}[h]
\centering
\begin{tabular}{p{0.34\linewidth}p{0.20\linewidth}p{0.38\linewidth}}
\toprule
\textbf{Item} & \textbf{Value} & \textbf{Evidence} \\
\midrule
Current branch & \texttt{main} & \texttt{git branch --show-current} \\
Commit count in current history & 17 & \texttt{git rev-list --count HEAD} \\
Tracked modified files & 2 & \path{.DS_Store}, \path{dataset/.DS_Store} \\
Untracked entries in working tree & 27 & \texttt{git status --short} \\
Experiment paths with git history & none & \texttt{git log -- <experiment path>} returns empty for Granite/Qwen/DeBERTa/RoBERTa folders \\
\bottomrule
\end{tabular}
\caption{Git provenance snapshot (February 14, 2026).}
\end{table}

\noindent\textbf{Interpretation.} The model experimentation work is primarily filesystem-level (untracked artifacts), not commit-level. This affects reproducibility and diff-based narrative writing.

\section{Artifact Footprint by Model Family}
\begin{longtable}{p{0.12\linewidth}p{0.10\linewidth}p{0.20\linewidth}p{0.52\linewidth}}
\toprule
\textbf{Family} & \textbf{File Count} & \textbf{Top Extensions} & \textbf{Representative Artifacts} \\
\midrule
\endfirsthead
\toprule
\textbf{Family} & \textbf{File Count} & \textbf{Top Extensions} & \textbf{Representative Artifacts} \\
\midrule
\endhead
Granite & 23 & \texttt{.json}, \texttt{.py}, \texttt{.ipynb}, \texttt{.md} &
\path{GraniteClarityEvaluation.ipynb}; \path{GraniteClarityEvaluation copy 2.ipynb}; \path{train_granite_rationale.py}; \path{granite_clarity_finetuned/checkpoint-60/trainer_state.json} \\

Qwen & 3 & \texttt{.py}, \texttt{.ipynb} &
\path{QwenTuning.ipynb}; \path{train_qwen.py}; \path{qwen_classifier.py} \\

DeBERTa & 3 & \texttt{.json} &
\path{deberta/model_evaluation_results_qevasion.json}; \path{deberta/model_evaluation_results copy.json}; \path{deberta/model_evaluation_results.json} \\

RoBERTa & 85 & \texttt{.json}, \texttt{.txt}, \texttt{.bin}, \texttt{.yaml}, \texttt{.csv} &
\path{roberta/roberta-ensemble-model-1/checkpoint-388/trainer_state.json}; \path{roberta/roberta-ensemble-model-2/checkpoint-388/trainer_state.json}; \path{roberta/train_clarity_ensemble.py}; \path{roberta/roberta_ensemble_exported/roberta-ensemble-model-1/model.onnx} \\

Ensemble & 88 & \texttt{.json}, \texttt{.yaml}, \texttt{.csv}, \texttt{.py} &
\path{clarity_ensemble/ensemble_roberta-base_20260201_134652/fold-1/logs/version_0/metrics.csv}; \path{ensemble_train_results.json}; \path{evaluate_ensemble.py} \\
\bottomrule
\caption{Artifact footprint (substring-based family mapping; families overlap).}
\end{longtable}

\section{Data and Split Evidence}
\begin{table}[h]
\centering
\begin{tabular}{p{0.50\linewidth}p{0.16\linewidth}p{0.28\linewidth}}
\toprule
\textbf{Dataset Evidence} & \textbf{Count} & \textbf{Notes} \\
\midrule
\path{dataset/QAEvasion.csv} & 3448 rows & Matches \texttt{train\_size=3448} in ensemble config files \\
\path{dataset/clarity_task_evaluation_dataset.csv} & 237 rows & Matches submission pickle length (237) \\
QEvasion test split (notebook loader output) & 308 rows & Reported by Granite notebook logs before balancing \\
Balanced Granite evaluation subset & 69 rows & 23 per class in notebook logs (\texttt{Direct Reply}, \texttt{Direct Non-Reply}, \texttt{Indirect}) \\
\bottomrule
\end{tabular}
\caption{Dataset and split counts recovered from local artifacts.}
\end{table}

\section{Configuration Snapshot by Family}
\begin{longtable}{p{0.12\linewidth}p{0.38\linewidth}p{0.44\linewidth}}
\toprule
\textbf{Family} & \textbf{Core Configuration} & \textbf{Source} \\
\midrule
\endfirsthead
\toprule
\textbf{Family} & \textbf{Core Configuration} & \textbf{Source} \\
\midrule
\endhead
Granite &
Base model \texttt{ibm-granite/granite-3.2-2b-instruct}; LoRA: \texttt{r=8}, \texttt{alpha=32}, \texttt{dropout=0.05}; target modules include \texttt{q\_proj/k\_proj/v\_proj/o\_proj} and MLP projections. &
\path{granite_clarity_finetuned/adapter_config.json} \\

Granite training checkpoint &
\texttt{global\_step=60}, \texttt{num\_train\_epochs=2}; training loss decreases from 1.6479 to 1.2488; no stored eval metrics in checkpoint state. &
\path{granite_clarity_finetuned/checkpoint-60/trainer_state.json} \\

Qwen &
Training defaults: model \texttt{unsloth/Qwen3-8B-bnb-4bit}, mode \texttt{classification\_head|generative}, \texttt{max\_epochs=3}, \texttt{batch\_size=4}, \texttt{lr=1e-4}; LoRA enabled in classifier implementation. &
\path{train_qwen.py}, \path{qwen_classifier.py} \\

DeBERTa &
QEvasion eval file uses \texttt{microsoft/deberta-v3-base}, threshold 0.5; legacy binary evaluation file references \texttt{lxyuan/distilbert-base-multilingual-cased-sentiments-student}. &
\path{deberta/model_evaluation_results_qevasion.json}, \path{deberta/model_evaluation_results copy.json} \\

RoBERTa / Clarity ensemble &
\texttt{roberta-base}, PEFT enabled, \texttt{lora\_r=16}, \texttt{lora\_alpha=32}, \texttt{dropout=0.2}, \texttt{label\_smoothing=0.1}, \texttt{lr=2e-5}, \texttt{epochs=5}, \texttt{num\_folds=3}. &
\path{clarity_ensemble/ensemble_roberta-base_*/config.json}, \path{clarity_ensemble/**/hparams.yaml} \\
\bottomrule
\caption{Recovered configuration-level information by model family.}
\end{longtable}

\section{Consolidated Quantitative Results}
\begin{longtable}{p{0.11\linewidth}p{0.26\linewidth}p{0.09\linewidth}p{0.07\linewidth}cccccp{0.18\linewidth}}
\toprule
\textbf{Family} & \textbf{Artifact} & \textbf{Split} & \textbf{n} & {\textbf{Acc}} & {\textbf{F1 (macro)}} & {\textbf{Prec}} & {\textbf{Recall}} & {\textbf{ROC-AUC}} & \textbf{Notes} \\
\midrule
\endfirsthead
\toprule
\textbf{Family} & \textbf{Artifact} & \textbf{Split} & \textbf{n} & {\textbf{Acc}} & {\textbf{F1 (macro)}} & {\textbf{Prec}} & {\textbf{Recall}} & {\textbf{ROC-AUC}} & \textbf{Notes} \\
\midrule
\endhead
Granite &
\path{GraniteClarityEvaluation copy 2.ipynb} (cell 12 output) &
Balanced test & 69 & 0.4638 & 0.4282 & {--} & {--} & {--} & Best Granite notebook metric found. \\

Granite &
\path{GraniteClarityEvaluation.ipynb} (cells 26--28 outputs) &
Balanced test & 69 & 0.3333 & 0.1667 & 0.1111 & 0.3333 & {--} & Degenerate behavior (all predictions as \texttt{Indirect}). \\

DeBERTa &
\path{deberta/model_evaluation_results_qevasion.json} &
QEvasion test & 308 & 0.4156 & 0.3878 & 0.2651 & 0.7215 & 0.5402 & 3-class clarity evaluation JSON. \\

DeBERTa (binary legacy) &
\path{deberta/model_evaluation_results copy.json} &
Binary eval & 200 & 0.7100 & 0.3830 & 0.5000 & 0.3103 & {--} & Binary setup, not directly comparable to 3-class runs. \\

RoBERTa fold run &
\path{clarity_ensemble/ensemble_roberta-base_20260201_134652/fold-1/logs/version_0/metrics.csv} &
Test & 308 & 0.4253 & 0.3522 & 0.3683 & 0.4314 & {--} & Best complete fold-level test line available. \\

RoBERTa fold run &
\path{clarity_ensemble/ensemble_roberta-base_20260201_134652/fold-1/logs/version_0/metrics.csv} &
Best val & {--} & 0.3913 & 0.3740 & 0.3985 & 0.4717 & {--} & Best logged validation epoch. \\
\bottomrule
\caption{Primary comparable metrics across model families from local files and notebook outputs.}
\end{longtable}

\subsection*{Legacy RoBERTa checkpoint diagnostics (not used in comparable ranking)}
\begin{table}[h]
\centering
\begin{tabular}{p{0.36\linewidth}cccp{0.30\linewidth}}
\toprule
\textbf{Checkpoint Artifact} & \textbf{Seed} & \textbf{Epochs} & \textbf{eval\_macro\_f1} & \textbf{Interpretation} \\
\midrule
\path{roberta/roberta-ensemble-model-1/checkpoint-388/trainer_state.json} & 42 & 2 & 1.0000 & Internal checkpoint eval; likely non-comparable to fold metrics. \\
\path{roberta/roberta-ensemble-model-2/checkpoint-388/trainer_state.json} & 123 & 2 & 1.0000 & Internal checkpoint eval; likely non-comparable to fold metrics. \\
\path{roberta/roberta-ensemble-model-3/checkpoint-388/trainer_state.json} & 456 & 1 & 0.1786 & Different training duration and behavior from model-1/2. \\
\bottomrule
\end{tabular}
\caption{Legacy HF Trainer checkpoint diagnostics. Seeds/epochs recovered from training\_args.bin artifacts.}
\end{table}

\section{Granite Notebook Deep-Dive}
\begin{table}[h]
\centering
\begin{tabular}{p{0.33\linewidth}ccccc}
\toprule
\textbf{Granite Notebook Run} & {\textbf{Overall Acc}} & {\textbf{Macro F1}} & {\textbf{F1 Direct Reply}} & {\textbf{F1 Direct Non-Reply}} & {\textbf{F1 Indirect}} \\
\midrule
\path{GraniteClarityEvaluation copy 2.ipynb} & 0.4638 & 0.4282 & 0.6538 & 0.2308 & 0.4000 \\
\path{GraniteClarityEvaluation.ipynb} & 0.3333 & 0.1667 & 0.0000 & 0.0000 & 0.5000 \\
\bottomrule
\end{tabular}
\caption{Granite evaluation metrics extracted from notebook output cells.}
\end{table}

\noindent Additional evidence from \path{GraniteClarityEvaluation.ipynb} confusion matrix:
\[
\begin{bmatrix}
0 & 0 & 23 \\
0 & 0 & 23 \\
0 & 0 & 23
\end{bmatrix}
\]
which confirms one-label collapse in that run.

\section{DeBERTa and Ensemble JSON Results}
\subsection{DeBERTa result files}
\begin{table}[h]
\centering
\begin{tabular}{p{0.43\linewidth}p{0.17\linewidth}p{0.32\linewidth}}
\toprule
\textbf{File} & \textbf{Examples} & \textbf{Key Signals} \\
\midrule
\path{deberta/model_evaluation_results_qevasion.json} & 308 &
3-class metrics present (Acc/F1/Prec/Recall/ROC-AUC). \\
\path{deberta/model_evaluation_results copy.json} & 200 &
Binary metrics with full prediction vectors; older/alternate configuration. \\
\path{deberta/model_evaluation_results.json} & 237 &
Submission-like prediction counts only, no quality metrics. \\
\bottomrule
\end{tabular}
\caption{DeBERTa result-file semantics.}
\end{table}

\subsection{Ensemble train summary file}
\begin{table}[h]
\centering
\begin{tabular}{p{0.27\linewidth}cccc}
\toprule
\textbf{Entry in \path{ensemble_train_results.json}} & {\textbf{Acc}} & {\textbf{Prec (macro)}} & {\textbf{Recall (macro)}} & {\textbf{F1 (macro)}} \\
\midrule
\texttt{model\_1\_train} & 0.5916 & 0.1972 & 0.3333 & 0.2478 \\
\texttt{model\_2\_train} & 0.5916 & 0.1972 & 0.3333 & 0.2478 \\
\texttt{model\_3\_train} & 0.3715 & 0.3074 & 0.3356 & 0.2558 \\
\texttt{ensemble\_average\_train} & 0.5916 & 0.1972 & 0.3333 & 0.2478 \\
\bottomrule
\end{tabular}
\caption{Ensemble summary JSON values (train split only).}
\end{table}

\section{RoBERTa/Ensemble Run-Health Audit}
\begin{longtable}{p{0.22\linewidth}p{0.10\linewidth}p{0.10\linewidth}p{0.14\linewidth}p{0.38\linewidth}}
\toprule
\textbf{Run Directory} & \textbf{Config} & \textbf{Metrics CSV} & \textbf{Fold Dirs} & \textbf{Status} \\
\midrule
\endfirsthead
\toprule
\textbf{Run Directory} & \textbf{Config} & \textbf{Metrics CSV} & \textbf{Fold Dirs} & \textbf{Status} \\
\midrule
\endhead
\texttt{20260201\_132359} & yes & yes & only \texttt{fold-1} & Full training/val log exists; no test line recorded in CSV. \\
\texttt{20260201\_134652} & yes & yes & only \texttt{fold-1} & Most complete run in this folder family; includes one test metric row. \\
\texttt{20260201\_135710} & yes & yes & only \texttt{fold-1} & Truncated (3 rows total). \\
\texttt{20260201\_185917} & yes & yes & only \texttt{fold-1} & Partial (8 rows total). \\
\texttt{20260202\_135909} & yes & no & only \texttt{fold-1} & Config + hparams exist; no metrics CSV written. \\
\bottomrule
\caption{Run completeness for the clarity\_ensemble roberta-base run directories.}
\end{longtable}

\noindent\textbf{Critical reproducibility mismatch.} Configs specify \texttt{num\_folds=3}, but every run stores only \texttt{fold-1} artifacts in this workspace snapshot.

\subsection*{Provenance of \texttt{ensemble\_20260201\_134652\_valbest}}
The identifier \texttt{ensemble\_20260201\_134652\_valbest} is a post-hoc report label (not a file name). It is defined as:
\[
\arg\max_{r \in R} \; r[\texttt{val/f1\_macro}]
\]
where \(R\) is the set of rows in:
\path{clarity_ensemble/ensemble_roberta-base_20260201_134652/fold-1/logs/version_0/metrics.csv}
with non-empty \texttt{val/f1\_macro}.

\begin{table}[h]
\centering
\begin{tabular}{ccccc}
\toprule
\textbf{Step} & \textbf{Epoch} & \textbf{val/acc} & \textbf{val/f1\_macro} & \textbf{val/loss} \\
\midrule
71  & 0 & 0.1035 & 0.0625 & 1.1725 \\
143 & 1 & 0.1200 & 0.0827 & 1.1637 \\
215 & 2 & 0.2730 & 0.2760 & 1.1507 \\
287 & 3 & 0.3548 & 0.3462 & 1.1405 \\
\textbf{359} & \textbf{4} & \textbf{0.3913} & \textbf{0.3740} & \textbf{1.1365} \\
\bottomrule
\end{tabular}
\caption{Validation checkpoints used to select \texttt{ensemble\_20260201\_134652\_valbest}.}
\end{table}

Pipeline-wise, this run is created by:
\begin{enumerate}
    \item Writing run config and timestamp (\path{roberta/train_clarity_ensemble.py}, lines 618--627).
    \item Building stratified folds (\path{roberta/train_clarity_ensemble.py}, lines 630--650).
    \item Training with checkpoint monitor \texttt{val/f1\_macro} (\path{roberta/train_clarity_ensemble.py}, lines 517--531).
    \item Logging metrics via Lightning CSV logger (\path{roberta/train_clarity_ensemble.py}, line 536), yielding the row sequence above.
\end{enumerate}

For the same run, the subsequent test row is logged at step 360, epoch 5, with \texttt{test/f1\_macro=0.3522}.

\section{Qwen Experiment Audit}
\begin{table}[h]
\centering
\begin{tabular}{p{0.32\linewidth}p{0.62\linewidth}}
\toprule
\textbf{Artifact} & \textbf{Observed Status} \\
\midrule
\path{QwenTuning.ipynb} & Recorded runtime failure: \texttt{torch\_shm\_manager ... Permission denied} (DataLoader worker). Follow-up \texttt{NameError: model is not defined}. \\
\path{Untitled3.ipynb} & Recorded \texttt{NameError: model is not defined}. \\
\path{train_qwen.py}, \path{qwen_classifier.py} & Training/evaluation code exists with LoRA paths and metric logging, but no saved checkpoint (\texttt{.ckpt}) files found in repository. \\
Qwen metrics artifacts & No standalone Qwen \texttt{.json}/\texttt{.csv} result file found for final numeric reporting. \\
\bottomrule
\end{tabular}
\caption{Qwen status: implementation present, but run outputs are incomplete in current workspace.}
\end{table}

\section{Submission Artifact Audit}
\begin{longtable}{p{0.38\linewidth}p{0.16\linewidth}p{0.40\linewidth}}
\toprule
\textbf{Artifact} & \textbf{Size/Count} & \textbf{Audit Notes} \\
\midrule
\endfirsthead
\toprule
\textbf{Artifact} & \textbf{Size/Count} & \textbf{Audit Notes} \\
\midrule
\endhead
\path{clarity/clarity_submission.pickle} & 237 predictions &
Tuple format valid (indices + labels). Label distribution: 126 \texttt{Clear Reply}, 111 \texttt{Non-Clear / Ambiguous Reply}. \\

\path{clarity/clarity_codebench_submission.csv} & 237 rows &
All rows have \texttt{prediction=2}, \texttt{label=2}. This appears degenerate and should not be treated as a final competitive submission without validation. \\

\path{generate_clarity_submission.py} & script &
Implements submission generation pipeline via \texttt{SimpleClarityClassifier}; output label mapping differs from 3-class CLARITY naming used in Granite notebooks. \\

\path{GraniteClarityEvaluation*.ipynb} & notebook logic &
Writes \path{clarity_submission_granite.pickle} in code cells, but file is not present in current workspace snapshot. \\
\bottomrule
\caption{Submission-related files and consistency checks.}
\end{longtable}

\section{Key Findings}
\begin{enumerate}
    \item \textbf{Best recoverable Granite metric:} \textbf{Acc 0.4638, Macro-F1 0.4282} from \path{GraniteClarityEvaluation copy 2.ipynb} on the balanced 69-sample subset.
    \item \textbf{Best complete RoBERTa fold test metric in \path{clarity_ensemble}:} \textbf{Acc 0.4253, Macro-F1 0.3522} from run \texttt{20260201\_134652}.
    \item \textbf{Best DeBERTa 3-class JSON metric:} \textbf{Acc 0.4156, Macro-F1 0.3878} on 308 examples.
    \item \textbf{Qwen runs are not yet reportable} due notebook/runtime failures and absence of saved checkpoints/results.
    \item \textbf{Git provenance is weak for experiments} because nearly all model artifacts are untracked; formal paper claims should cite filesystem evidence and dates explicitly.
\end{enumerate}

\section{Recommended Finalization Plan for SemEval Submission Write-Up}
\begin{enumerate}
    \item Promote the complete experiment artifacts (at minimum: configs, metrics CSV/JSON, and selected notebooks) into tracked git commits so paper tables are reproducible from commit hashes.
    \item Normalize one canonical label space (\texttt{Direct Reply}, \texttt{Direct Non-Reply}, \texttt{Indirect}) across Granite/DeBERTa/RoBERTa/legacy binary artifacts before final leaderboard submission.
    \item Re-run Qwen with fixed DataLoader settings (\texttt{num\_workers=0} and verified shared-memory permissions), then export a dedicated metrics JSON comparable with other families.
    \item Complete true 3-fold outputs for \path{clarity_ensemble} or explicitly document single-fold limitation.
    \item Generate one final, validated submission file from the selected best model and store it with a dated manifest (model, commit hash, metrics source paths).
\end{enumerate}

\end{document}
